% ---------------------------------------------------------------------
%  APÊNDICE D — ATITUDE VERDADEIRA E ATITUDE DO AHRS
%  Revisão de 3/11, item 2.2 (P3c): as tabelas usam a atitude verdadeira;
%  aqui, as duas bases nos mesmos voos. Tabela gerada por
%  ensaios/tabelas_monografia.py (apendice_bases.tex).
% ---------------------------------------------------------------------
\chapter{Atitude verdadeira e atitude do AHRS}
\label{ap:bases}

As tabelas de resultados medem a atitude verdadeira do modelo físico do SITL
(mensagem \codigo{CCTQ}), menos o viés médio entre as bases no meio segundo
antes da perturbação (Seção~\ref{sec:metricas}). A Tabela~\ref{tab:bases} mostra, nos mesmos
voos, o desvio máximo e o erro eficaz nas duas bases, para o nativo e os dois
LQG nos ensaios em que a diferença mais pesa: o impulso, o degrau e o torque
estocástico, sem ruído e com o ruído de giroscópio de
\codigo{SIM\_GYR*\_RND} = 20.

\pendente{Fase 6: comentar a tabela com os voos do lote (onde as bases
divergem, e se a ordem entre os controladores muda).}

\begin{table}[htbp]
\centering
\caption{Desvio máximo e erro eficaz na atitude do AHRS e na atitude
verdadeira, nos mesmos voos.}
\label{tab:bases}
\resizebox{\textwidth}{!}{% Gerada por ensaios/tabelas_monografia.py — nao editar a mao.
\begin{tabular}{lllrrrrrr}
\toprule
 & & & \multicolumn{3}{c}{Desvio máx. (\si{\degree})} & \multicolumn{3}{c}{RMSE (\si{\degree})} \\
\cmidrule(lr){4-6}\cmidrule(lr){7-9}
Ensaio & Eixo & Controlador & AHRS & verdade & $\Delta$ (\%) & AHRS & verdade & $\Delta$ (\%) \\
\midrule
impulso de rolagem & Rolagem & nativo & 0,789 & 0,802 & +2 & 0,080 & 0,081 & +2 \\
impulso de rolagem & Rolagem & LQG & 0,731 & 0,744 & +2 & 0,088 & 0,091 & +4 \\
impulso de rolagem & Rolagem & LQG (mesmo ruído) & --- & --- & --- & --- & --- & --- \\
\addlinespace
degrau de rolagem & Rolagem & nativo & 0,599 & 0,614 & +2 & 0,070 & 0,086 & +23 \\
degrau de rolagem & Rolagem & LQG & 0,597 & 0,612 & +3 & 0,039 & 0,059 & +52 \\
degrau de rolagem & Rolagem & LQG (mesmo ruído) & --- & --- & --- & --- & --- & --- \\
\addlinespace
estocástico de rolagem & Rolagem & nativo & --- & --- & --- & --- & --- & --- \\
estocástico de rolagem & Rolagem & LQG & --- & --- & --- & --- & --- & --- \\
estocástico de rolagem & Rolagem & LQG (mesmo ruído) & --- & --- & --- & --- & --- & --- \\
\addlinespace
impulso de guinada & Guinada & nativo & --- & --- & --- & --- & --- & --- \\
impulso de guinada & Guinada & LQG & --- & --- & --- & --- & --- & --- \\
impulso de guinada & Guinada & LQG (mesmo ruído) & --- & --- & --- & --- & --- & --- \\
\addlinespace
degrau de guinada & Guinada & nativo & --- & --- & --- & --- & --- & --- \\
degrau de guinada & Guinada & LQG & --- & --- & --- & --- & --- & --- \\
degrau de guinada & Guinada & LQG (mesmo ruído) & --- & --- & --- & --- & --- & --- \\
\addlinespace
estocástico de guinada & Guinada & nativo & --- & --- & --- & --- & --- & --- \\
estocástico de guinada & Guinada & LQG & --- & --- & --- & --- & --- & --- \\
estocástico de guinada & Guinada & LQG (mesmo ruído) & --- & --- & --- & --- & --- & --- \\
\addlinespace
impulso (ruido/rnd20) & Rolagem & nativo & --- & --- & --- & --- & --- & --- \\
impulso (ruido/rnd20) & Rolagem & LQG & --- & --- & --- & --- & --- & --- \\
impulso (ruido/rnd20) & Rolagem & LQG (mesmo ruído) & --- & --- & --- & --- & --- & --- \\
\addlinespace
impulso (ruido/rnd20) & Guinada & nativo & --- & --- & --- & --- & --- & --- \\
impulso (ruido/rnd20) & Guinada & LQG & --- & --- & --- & --- & --- & --- \\
impulso (ruido/rnd20) & Guinada & LQG (mesmo ruído) & --- & --- & --- & --- & --- & --- \\
\bottomrule
\end{tabular}
}

\fonte{Elaborada pelo autor, por \codigo{ensaios/tabelas\_monografia.py}.}
\nota{\raggedright $\Delta$: verdade em relação ao AHRS.
% Gerada por ensaios/tabelas_monografia.py — nao editar a mao.
Os mesmos voos nas duas bases; um voo por célula, exceto: média de 3 voos para nativo, LQG, LQG (mesmo ruído) em degrau de rolagem, impulso de rolagem; média de 5 voos para nativo, LQG, LQG (mesmo ruído) em estocástico de guinada, estocástico de rolagem, impulso (ruido/rnd20).
}
\end{table}
